\section{Distance to singular pencils}\label{sec:distance}

\begin{proposition}\label{prop:distance}
Fix $R_1,\ldots,R_d\in M_m(\mathbb C)$ and $M>0$.
For a tuple $H=(H_1,\ldots,H_d)$ of Hermitian $N\times N$ matrices, use the pencil $T_N(H)$ from \eqref{eq:pencil}.
There exist $C,\alpha,\varepsilon_0>0$ such that, for every $N$ and every such $H$ with $\max_j\|H_j\|\le M$, the following holds.
If $s_{\min}(T_N(H))\le\varepsilon\le\varepsilon_0$ with $\varepsilon>0$, there are Hermitian $N\times N$ matrices $\Delta H_1,\ldots,\Delta H_d$ such that $T_N(H+\Delta H)$ is singular and
\begin{equation}\label{eq:perturbation-bound}
 \left(\sum_{j=1}^d\|\Delta H_j\|_{\mathrm{HS}}^2\right)^{1/2}
 \le C\varepsilon^\alpha.
\end{equation}
Here $\|X\|_{\mathrm{HS}}:=(\Tr(X^*X))^{1/2}$ is the Hilbert--Schmidt norm, with the unnormalized trace on the appropriate matrix size.
\end{proposition}

\begin{proof}
Choose a unit vector $v=(v_1,\ldots,v_m)\in\mathbb C^m\otimes\mathbb C^N$ with $\|T_N(H)v\|\le\varepsilon$.
Let $W:=\operatorname{span}\{v_1,\ldots,v_m\}$, and let $W_1:=\sum_{j=1}^d P_{W^\perp}H_jW$, where $P_{W^\perp}$ is the orthogonal projection onto $W^\perp$.
Writing $k:=\dim W$ and $\ell:=\dim W_1$, we have $1\le k\le m$ and $0\le\ell\le dk$.
Choose isometries $U:\mathbb C^k\to\mathbb C^N$ and $V:\mathbb C^\ell\to\mathbb C^N$ onto $W$ and $W_1$, respectively, with $\mathbb C^0=\{0\}$.
For each $j$ put $A_j:=U^*H_jU$ and $B_j:=V^*H_jU$, so that $A_j$ is Hermitian and $B_j\in M_{\ell,k}(\mathbb C)$.
By the definition of $W_1$,
\begin{equation}\label{eq:columns-in-span}
 H_jU=UA_j+VB_j.
\end{equation}
Write $w:=(U^*v_1,\ldots,U^*v_m)\in\mathbb C^m\otimes\mathbb C^k$.
The two coordinate vectors of $T_N(H)v$ in the orthogonal summands $\mathbb C^m\otimes W$ and $\mathbb C^m\otimes W_1$ are
\begin{equation}\label{eq:compressed-residual}
 w+\sum_{j=1}^d(R_j\otimes A_j)w
 \quad\text{and}\quad
 \sum_{j=1}^d(R_j\otimes B_j)w.
\end{equation}
Equation~\eqref{eq:columns-in-span} gives no component outside these summands, so the sum of the squared norms in \eqref{eq:compressed-residual} is $\|T_N(H)v\|^2$.

Use Hilbert--Schmidt orthonormal real coordinates for each $A_j$, and split the entries of each $B_j$ and of $w$ into real and imaginary parts.
For each $(k,\ell)$ in the indicated ranges, these choices give $p_{k,\ell}:=dk^2+2d\ell k+2mk$ real coordinates for $x:=(A,B,w)$.
Let $F_{k,\ell}:\mathbb R^{p_{k,\ell}}\to\mathbb R^{2m(k+\ell)+1}$ be the polynomial map whose coordinates are the real and imaginary parts of \eqref{eq:compressed-residual} and the scalar $\|w\|^2-1$.
The map is quadratic and depends only on the coefficients $R_j$ and the pair $(k,\ell)$.
Since $\|w\|=1$ and $\|A_j\|_{\mathrm{HS}},\|B_j\|_{\mathrm{HS}}\le\sqrt{k}M$, the Euclidean norms of $x$ and its image satisfy
\begin{equation}\label{eq:compressed-input-bounds}
 |x|^2\le 2dkM^2+1,
 \qquad
 |F_{k,\ell}(x)|\le\varepsilon.
\end{equation}

The real polynomial distance inequality of Kurdyka and Spodzieja, as stated in \cite[p.~2, \textup{(KS2)}]{KSS16}, gives the following bound for a polynomial map $F:\mathbb R^p\to\mathbb R^q$ of degree $D\ge1$ with a nonempty zero set.
For some $c_F>0$ and every $y\in\mathbb R^p$,
\[
 |F(y)|\ge c_F
 \left(\frac{\operatorname{dist}(y,F^{-1}(0))}{1+|y|^2}\right)^{D(6D-3)^{p-1}}.
\]
When $F_{k,\ell}$ has a zero, its zero set is closed, so the distance from $x$ is attained at some $x':=(A',B',w')$.
Using $D=2$ and \eqref{eq:compressed-input-bounds} gives
\begin{equation}\label{eq:compressed-repair}
 |x'-x|\le C_{k,\ell}\varepsilon^{\alpha_{k,\ell}},
 \qquad
 \alpha_{k,\ell}:=(2\cdot 9^{p_{k,\ell}-1})^{-1},
\end{equation}
where $C_{k,\ell}$ depends only on the map and $M$.
If the zero set is empty, then $|F_{k,\ell}|$ has a positive minimum on the closed ball $|x|^2\le 2dkM^2+1$.
There are only finitely many pairs $(k,\ell)$, so we can choose $\varepsilon_0\le1$ below these positive minima, and choose a common exponent $\alpha>0$ and a constant $C_1$ in \eqref{eq:compressed-repair}.

With the original $U$ and $V$, define
\[
 \Delta H_j:=U(A'_j-A_j)U^*+V(B'_j-B_j)U^*+U(B'_j-B_j)^*V^*.
\]
This matrix is Hermitian and satisfies $(H_j+\Delta H_j)U=UA'_j+VB'_j$ by \eqref{eq:columns-in-span} and $V^*U=0$.
The coordinates of $T_N(H+\Delta H)(I_m\otimes U)w'$ are therefore the expressions in \eqref{eq:compressed-residual} with $A_j,B_j,w$ replaced by $A'_j,B'_j,w'$.
Both expressions vanish because $F_{k,\ell}(x')=0$, and its normalization coordinate gives $\|(I_m\otimes U)w'\|=1$.
This produces a kernel vector even if its components span a proper subspace of $W$.
The three blocks in $\Delta H_j$ are pairwise Hilbert--Schmidt orthogonal, so
\[
 \sum_{j=1}^d\|\Delta H_j\|_{\mathrm{HS}}^2
 =\sum_{j=1}^d\big(\|A'_j-A_j\|_{\mathrm{HS}}^2
                  +2\|B'_j-B_j\|_{\mathrm{HS}}^2\big)
 \le 2|x'-x|^2.
\]
This proves \eqref{eq:perturbation-bound} with $C=\sqrt{2}C_1$.
\end{proof}
